% !TeX root = main.tex
\documentclass[main.tex]{subfiles}

\begin{document}


\noindent
In this investigation, we evaluate the effectiveness of a custom BERT-based model in static malware detection. The model is trained to distinguish between malicious and benign opcode sequences using modern natural language processing techniques such as pre-training, which are then used to identify malicious programs. We compare this model against both linear models, such as decision trees, and other static machine learning methods seen in the literature, such as Convolutional Neural Networks. Finally, we evaluate to what extent the opcode sequence classification, and the CNN image classification models rely on the presence of obfuscation as an indicator of malicious software.

\section{Dataset}\label{sec:dataset}

\begin{figure}
    \centering
    \begin{subfigure}{0.4\textwidth}
        \includegraphics[width=\textwidth]{size-dist.png}
        \caption{Distribution of executable size across classes.}
        \label{fig:dataset:size-dist}
    \end{subfigure}
    \hfill
    \begin{subfigure}{0.4\textwidth}
        \includegraphics[width=\textwidth]{obfuscation-dist.png}
        \caption{Distribution of code analysis coverage across classes.}
        \label{fig:dataset:obf-dist}
    \end{subfigure}
    % \hfill
    \vspace{1em}
    \\
    \centering
    \begin{subfigure}{0.4\textwidth}
        \centering
        \caption{Distribution of classes with our dataset.}
        \begin{tabular}{@{}ll@{}}
            \toprule
            \textbf{Class}     & \textbf{Count} \\ \midrule
            Malicious & 5792   \\
            Benign    & 5323   \\ \midrule
            \textbf{Total}     & 11115  \\ \bottomrule
        \end{tabular}
        \label{fig:dataset:class-dist}
    \end{subfigure}
    \caption{Description of our dataset, including the size of programs (a), the extent of obfuscation (b), and the total number of samples in each class (c).}
    \label{fig:dataset}
\end{figure}

Our dataset consists of 11115 Windows PE32 files (See \cref{apdx:pe}), with a class distribution shown in \cref{fig:dataset:class-dist}. Malicious files are gathered from VirusShare\footnote{\url{https://virusshare.com/about}}, which contains samples from real-world cybersecurity incidents, providing approximately 90 million samples of which nearly 6000 were randomly selected. The benign executables were collected from the personal devices of the authors, and were scanned using Microsoft defender to ensure no mislabelled samples were present. This includes various software packages, games, developer tools, system programs, and other miscellaneous executables to ensure the dataset is representative of the kinds of benign programs commonly seen on the machines of potential end users.

We opt to investigate Windows malware detection due to the operating systems widespread usage, and therefore high availability of malware targeting it. Malware targeting other operating systems, such as Linux of MacOS do exist, but there are significantly fewer samples, meaning the availability of data is reduced.

The dataset demonstrates wide variation in characteristics between classes. For example, we find that malicious programs are often much smaller than benign executables, as shown in \cref{fig:dataset:size-dist}. We believe this to be because benign programs typically provide more functionality, such as graphical user interfaces, telemetry, and compatibility, which increase the size of the program. Malicious software on the other hand, are developed with a single goal in mind, reducing its size.

We hypothesize that both our dataset, and other datasets containing malicious software will demonstrate variation in the prevalence of obfuscation between malicious and benign executable classes. To investigate this, we perform static analysis of the samples in our dataset using Rizin \citep{rizinorganizationRizin2025} to find the code analysis coverage. This metric gives the ratio of code analysed by the tool, to the total amount of code in an executable. In obfuscated programs, this result is typically much lower than non-obfuscated executables as it is unable to perform its analysis due to the non-standard structure. The results of this analysis are presented in \cref{fig:dataset:obf-dist}, which supports our hypothesis, and indicates malicious samples are more frequently obfuscated than benign samples.  

The dataset is then split to produce a training set, and evaluation set at a ratio of 80\% to 20\% of samples respectively. To ensure consistent class distribution between these sets, samples are shuffled prior to splitting. 

\subsection{Obfuscation Experiment Dataset}\label{ssec:secondary-dataset}

\begin{table}[h]
    \centering
    \caption{Distribution of classes with our secondary dataset.}
    \begin{tabular}{@{}ll@{}}
        \toprule
        \textbf{Class}      & \textbf{Count} \\ \midrule
        Malicious           & 60   \\
        Benign              & 61   \\ \midrule
        \textbf{Total}      & 121  \\ \bottomrule
    \end{tabular}
    \label{fig:obfs-ds-class-dist}
\end{table}

\noindent
To investigate the impact of obfuscation on malware detection, we also curate a secondary dataset, consisting of unobfuscated executables. We use the work presented in \citep{rokonSourceFinderFindingMalware2020} to identify GitHub repositories containing malicious source code. The authors provide a list of URLs to these repositories, from which we select Windows programs that provide an executable and with no mention of obfuscation in its description or codebase. Benign samples consist of windows system binaries not present in our original dataset. We identify a collection of approximately 60 Windows XP system binaries which are free of obfuscation for this secondary dataset. This gives a selection of 121 programs, each of which are free of obfuscation for use in our later experiment into the resilience of machine learning algorithms to obfuscation. 

\section{Metrics}

\begin{figure}[H]
    \centering
    \includegraphics{confusion-matrix.pdf}
    \caption{Confusion Matrix describing the kinds of classification results in malware detection.}
    \label{fig:confusion-matrix}
\end{figure}

In our investigation, malware detection is presented as a binary classification problem, where our models aim to distinguish between malicious and benign software. To quantify the performance of the various approaches tested in this investigation, we find the accuracy, F1-score, precision, and recall~\cref{eq:metrics}. These metrics are based on the different kinds of results a binary classifier can produce, shown in \cref{fig:confusion-matrix}. 

Accuracy is the ratio of correct classifications to the number of samples tested, where a score of ``1'' would be a perfect classifier for a given set of samples. However, the value of this metric can be misleading in unbalanced datasets. 

Precision and recall both operate on the class-level, therefore to reduce these metrics to a single number, we typically take the average of these values for each class in the dataset, known as macro-averaging. Precision is defined as the ratio between the number of correctly identified samples, to the total number of samples predicted as belonging to a given class. Recall on the other hand is the number of correct predictions, to the number of samples in given class. F1 score, acts as an aggregation of these values, and is given by the harmonic mean.

\begin{align}
    \text{Accuracy} &= \frac{TP + TN}{TP + TN + FP + FN}\nonumber\\
    \text{Precision} &= \frac{TP}{TP + FP} \refstepcounter{equation} \tag{\theequation}\label{eq:metrics}\\
    \text{Recall} &= \frac{TP}{TP + FN}\nonumber\\ 
    \text{F1-Score} &=  \frac{2 \cdot \text{Precision} \cdot \text{Recall}}{\text{Precision} + \text{Recall}}\nonumber
\end{align}


\section{Baseline}

\begin{figure}[h]
    \begin{subfigure}{0.4\textwidth}
        \centering
        \begin{tabular}{@{}l|rr@{}}
        \toprule
        \textbf{Similarity} & \multicolumn{1}{l}{Malicious} & \multicolumn{1}{l}{Benign} \\ \midrule
        Malicious           & 0.1428                        & 0.0931                     \\
        Benign              & 0.0931                        & 0.1217                     \\ \bottomrule
        \end{tabular}
        \caption{Cross-class external function similarity.}
    \end{subfigure}
    \hfill
    \begin{subfigure}{0.4\textwidth}
        \centering
        \begin{tabular}{@{}l|rr@{}}
        \toprule
        \textbf{Similarity} & \multicolumn{1}{l}{Malicious} & \multicolumn{1}{l}{Benign} \\ \midrule
        Malicious           & 0.2249                        & 0.1606                     \\
        Benign              & 0.1606                        & 0.1643                     \\ \bottomrule
        \end{tabular}
        \caption{Cross-class DLL similarity.}
    \end{subfigure}
    \caption{Cosine similarity between classes based on external functions and DLLs.}
    \label{fig:similarity}
\end{figure}

\noindent 
We firstly establish a baseline using conventional machine learning algorithms on the lookup table associated with each program. This lookup table tells the program where external references, such as libraries and functions exist in memory. These external references offer functionality that would otherwise be inaccessible to the program. For example, to interact with the file-system, the lookup table will include a reference to \texttt{KERNEL32.dll}, which provides the address of the functions needed to create, read, and write to the file-system. We can extract all functions and libraries used by an executable with other static analysis tools such as Radare2~\citep{radareorganizationRadare22025}. References to functions that occur once across the dataset are removed as approximately half of all functions are seen exactly once across the entire dataset. This gives two binary feature vectors for each sample in the dataset, telling the classifiers which external functions, and libraries are used by the program. We test 3 machine learning algorithms to establish this baseline, which are support vector classifiers, logistic regression, and decision trees. 

We also use the feature vectors to quantify the similarity of programs in our dataset with the pairwise cosine similarity function (\cref{eq:cosine-sim}). We then find the average similarity between classes which are presented in \cref{fig:similarity}. From this, we can assume that functions and libraries used by the malicious programs are more similar than the ones used by benign programs. Furthermore, functions used by malicious programs are vastly different to functions used by benign programs. 

\begin{equation}
    S\left(\mathbf{A},\mathbf{B}\right) := \cos\left(\theta\right)=\frac{\mathbf {A} \cdot \mathbf {B}}{\|\mathbf {A} \|\|\mathbf {B} \|}\label{eq:cosine-sim}
\end{equation}

\section{RoBERTa for Opcode Sequence Classification}

\begin{figure}[h]
    \centering
    \begin{BVerbatim}
ret lea mov add add lea mov add add nop
    \end{BVerbatim}
    \caption{An example sequence of opcodes}
    \label{fig:opcode-sequence}
\end{figure}

\noindent
In recent years, transformer-based architectures for natural language processing have been shown to be highly effective across a wide variety of language based tasks. A significant advancement within the field was the development of the Bidirectional Encoded Representation from Transformers model, known as BERT~\citep{devlinBERTPretrainingDeep2019a}. This model uses an encoder-only architecture, based on the multi-head self-attention mechanism~\citep{vaswaniAttentionAllYou2023a} in conjunction with a pre-training stage to develop a comprehensive understanding of language. Pre-training consists of two tasks, masked language modelling where the model predicts hidden tokens based on their surrounding context, and next sentence prediction. The pre-trained model can then be fine-tuned on a downstream problem such as sequence classification, question answering, and general language understanding. At the time of publication, it achieved state-of-the-art performance in the General Language Understanding Evaluation (GLUE), Stanford Question Answering Dataset (SQuAD), and Situations With Adversarial Generations (SWAG) benchmarks and has been used extensively throughout the natural language processing field. Fine-tuning of the pre-trained model allows researchers to solve complex tasks with fewer training samples, short training periods, and less intensive compute requirements. RoBERTa extends upon the original BERT model by expanding the scale of the model, removing the next sentence prediction task from its pre-training stage, optimising the hyperparameters, and training on a larger corpus. 

We develop a custom RoBERTa model, trained on opcode sequences for malware detection. Opcodes refer to the operation to be performed by the CPU, such as \texttt{add}, \texttt{cmp}, \texttt{mov}, written in a human readable form. To generate these sequences, we use Radare2 to extract the disassembly of a program (See \cref{algo:disasm-extraction}), and remove the operands of each instruction (\cref{algo:opcode-sequence-extraction}) to reduce the vocabulary size and therefore the computational resources needed to train the model. This gives a sequence of opcodes representing each program within our dataset, as shown in \cref{fig:opcode-sequence}. The sequences are then tokenized at the word level where each opcode is assigned a single token. 

We train four of these models that each take a tokenized opcode sequence of a fixed length. The first model is trained in masked language modelling, which is then used to train a sequence classification model, referred to as the ``pretrained'' model throughout this report. We also train another pre-trained sequence classification model which has had the weights in its attention layers frozen, referred to as the ``frozen'' model to determine if the representations produced by the attention layers are sufficient for distinguishing between malicious and benign opcode sequences. Finally, we also train a RoBERTa sequence classification model from a random initialisation, to assess if the pre-training step is necessary, referred to as the ``base'' model.  

Our models share the architecture shown in \cref{tab:model-params} that was found to be the optimal configuration after an extensive hyperparameter sweep of 90 runs. Training was performed with a single Nvidia A40 GPU, and 16GB of ram with a batch size of 128 for a single epoch, using the AdamW optimizer, and Cross-Entropy loss function provided by Pytorch.

For each opcode sequence, the model predicts two values, the likelihood of the sequence being malicious, or benign where the larger of these two values is its predicted class. However, the full opcode sequence of a program is too long to be used as the models input. Therefore, we split the full sequence of opcodes into multiple sub-sequences of a fixed length, and aggregate the logits to classify a program. We test both raw logit mean aggregation, softmax logit mean aggregation, and majority voting, and found raw logit mean aggregation to be most effective. 

\begin{table}[h]
    \centering
    \caption{RoBERTa sequence classification model parameters}
    \begin{tabular}{@{}lr@{}}
        \toprule
        \textbf{Parameter}      & \textbf{Value} \\ \midrule
        Sequence Length         & 32    \\
        Hidden Layer Size       & 512   \\
        Hidden Layers           & 7     \\
        Attention Heads         & 4     \\
        Intermediate Layer Size & 2048  \\
        Activation Function     & GELU  \\
        Hidden Dropout Rate     & 0.01  \\
        Attention Dropout Rate  & 0.001 \\ \bottomrule
    \end{tabular}
    \label{tab:model-params}
\end{table}

\section{CNN for Image Classification}

\begin{figure}
    \centering 
    \begin{subfigure}[t]{0.45\textwidth}
        \centering 
        \includegraphics[width=\textwidth]{malware-img-benign.png} 
        \caption{Image representation of a benign executable.}
    \end{subfigure}
    \hfill
    \begin{subfigure}[t]{0.45\textwidth}
        \centering 
        \includegraphics[width=\textwidth]{malware-img-mal.png} 
        \caption{Image representation of a malicious executable.}
    \end{subfigure}
    \caption{How images can be used to represent programs.}
    \label{fig:img-rep}
\end{figure}

In addition to the RoBERTa sequence classification model, we also evaluate an image based approach. We use the EfficientNet-B0 CNN architecture \citep{tanEfficientNetRethinkingModel2020}, due to its computational efficiency and high performance on image classification problems. The model is trained from a random initialisation, and the final layer of the model is modified to output a binary class likelihood values. Programs are converted to greyscale images using the procedure outlined in \cref{apdx:prog-to-img}, converted to RGB, and resized to 320\(\times\)320 pixels before it is used by the model, as shown in \cref{fig:img-rep}.

As previously established, the size of benign programs is typically much higher than malicious programs. Therefore, benign images are resized to a greater extent to that of malicious images. Resizing an image will introduce artifacts to the image, such as anti-aliasing or blurring. Thus, benign images may have more prevalent artifacts than malicious images which could be used by the model to learn a suboptimal solution. To prevent this from occurring, we train the same model on a subset of the dataset where the distribution of program size between classes are more similar, known as the ``filtered'' model. We use a genetic algorithm that minimises the Kullback-Leibler divergence $D_{KL}$ between the distributions of program size for each class of program by modelling it as a Knapsack problem~\citep{khuriZeroOneMultiple1994}. We also encourage the algorithm to select a subset with an approximately equal number of malicious and benign samples to ensure class balance. 

We train both the base model, and the filtered model trained for 20 epochs, with the binary cross entropy loss function and a batch size of 32. Training was conducted using an Apple M2 Pro CPU, leveraging the metal performance shaders backend for hardware acceleration, with 16GB of RAM using the AdamW optimiser. 

\section{Obfuscation}

In modern malware, obfuscation is highly pervasive. As a result of this, we believe that static malware detection methods learn to identify obfuscation as opposed to other malicious characteristics. However, obfuscation is not exclusive to malicious programs, as legitimate software developers use similar techniques to protect their intellectual property. Should this be the case, legitimate benign software that have been obfuscated may be misidentified as malicious, thereby degrading the effectiveness of these approaches in real world deployment. 

We investigate this by observing how the models performance changes as we increase the presence of obfuscation within the unobfuscated dataset. We perform three trials for each model, that vary which kinds of program are to be obfuscated. For example, only malicious, benign, and all programs. At each step, we obfuscate a certain percentage of programs within the dataset and this percentage increases by 10\% upon completion. We collect the logits for each sequence, which are used to calculate both the sequence level, and the file level metrics. We also perform this experiment on the image classification model to see if these models are better suited to obfuscation resilient malware detection.

Obfuscation is performed using a tool known as UPX~\citep{markusf.x.j.oberhumerUPX2025}, which is a binary packing tool. We chose this tool due to its extensive compatibility with Windows executables. Furthermore, a number of our malicious executables have been identified as using this tool for obfuscation, and it has also been used to compress and obfuscate benign files. Therefore, we feel that it is a suitable tool to replicate the forms of obfuscation used by both malicious and legitimate software developers seen in the wild without having to implement our own obfuscation tooling.

\end{document}